\subsection{Tauschoperationen im Spatprodukt}
\label{subsec:vector_algebra_scalar_triple_product_properties}

Das Spatprodukt bleibt bei bestimmten Änderungen der Reihenfolge unverändert,
während andere lediglich sein Vorzeichen ändern. Um beide Fälle sicher zu
unterscheiden, betrachten wir zunächst zyklische und paarweise Vertauschungen.

\begin{DefinitionBox}{Zyklische vs. Paarweise Vertauschung}{vector_algebra_cyclic_permutation}
  In einer geordneten Folge unterscheiden wir im Folgenden zwei Arten von
  Vertauschungen:
  \begin{itemize}[leftmargin=1.6em,itemsep=0.5em,topsep=0.3em]
    \item Bei einer \emph{zyklischen Vertauschung}
      \EnglishTerm{cyclic permutation} werden alle Glieder um eine Position
      verschoben. Das erste Glied rückt dabei an das Ende. Für drei Glieder gilt:
      \[
        (a,b,c)
        \longmapsto\Cloze{2.7cm}{(b,c,a)}
        \longmapsto\Cloze{2.7cm}{(c,a,b)}
        \longmapsto\Cloze{2.7cm}{(a,b,c)}.
      \]
    \item Eine \emph{paarweise Vertauschung}, kurz \emph{Paartausch}
      \EnglishTerm{transposition}, vertauscht genau zwei Glieder. Alle übrigen
      Glieder behalten ihre Position:
      \[
        (a,b,c)\longmapsto\Cloze{2.7cm}{(b,a,c)}.
      \]
  \end{itemize}
  Ein Paartausch ist daher im Allgemeinen keine zyklische Vertauschung.
\end{DefinitionBox}


Für das Spatprodukt zeigt der folgende Satz, wie sich beide Arten der
Vertauschung auf den Wert auswirken.
\begin{TheoremBox}{Vertauschungen im Spatprodukt}{vector_algebra_scalar_triple_product_cyclic_permutation}
  Das Spatprodukt $S=(\vec a\times\vec b)\cdot\vec c$
  \smallskip
  \begin{itemize}[leftmargin=1.6em,itemsep=0.2em,topsep=0pt]
    \item bleibt bei \textcolor{THMBlue}{zyklischer Vertauschung} der Vektoren
      unverändert,
    \item wechselt bei \textcolor{olive!80!black}{paarweiser Vertauschung} der
      Vektoren sein Vorzeichen.
  \end{itemize}
  \begin{center}
    \CachedTikzFigure{content/resources/chapters/vector_algebra/tikz/scalar_triple_product_permutations.tex}
  \end{center}
  Beide Vertauschungen lassen das Volumen unverändert, da $V=|S|$ gilt.
\end{TheoremBox}

\begin{ExampleBox}{Vertauschungen im Spatprodukt}{vector_algebra_scalar_triple_product_permutation_example}
  Für die Vektoren
  \[
    \vec a=\begin{pmatrix}3\\-1\\2\end{pmatrix},\qquad
    \vec b=\begin{pmatrix}4\\2\\1\end{pmatrix},\qquad
    \vec c=\begin{pmatrix}0\\2\\-5\end{pmatrix}
  \]
  berechnen und vergleichen wir drei Spatprodukte, die sich jeweils nur durch
  eine Vertauschung der Vektoren unterscheiden.
  \begin{center}
    \CachedTikzFigure{content/resources/chapters/vector_algebra/tikz/scalar_triple_product_permutation_example.tex}
  \end{center}
\end{ExampleBox}
